\chapter{The Shell}\label{cha:shell}

The shell is what the machine boots into. It is a file manager, a program
launcher and a front door to the monitor, and every command here also works
from BASIC with an \verb|@| in front.

\section{Files and directories}
Files live on the host (the \texttt{fs/} directory on the desktop, the
\texttt{k4510/fs/} folder of the SD card on the Pi). Programs are kept in
\texttt{/PRG}, BASIC text in \texttt{/BASIC}, music in \texttt{/SID} --- and a
bare name is searched there too, so you rarely need a path.

\begin{type}
DIR              CD BASIC         CD ..
MKDIR MYDIR      RM OLD.TXT       RMDIR MYDIR
TYPE README.BAS  LOAD name [addr] SAVE name from.to
\end{type}

\screen{shots/dir.png}{\texttt{DIR}: directories first in white, then files
with their sizes, two columns.}

\section{Running things}
\verb|RUN balls.prg| runs a program; so does \verb|RUN balls|, and so does
plain \verb|BALLS| --- an unknown word is tried as a program name before the
shell gives up. Try \verb|SIDPLAY|, the SID jukebox, or \verb|RUN ehbasic|
for BASIC.

\section{The monitor}
\verb|MON| opens the machine monitor at a \verb|*| prompt; its grammar is
Wozmon's, with 28-bit addresses:

\begin{type}
*E000.E00F           examine a range
*0300:A9 41 60       store bytes
*0300R               run from $0300
*X                   leave
\end{type}

\screen{shots/mon.png}{The monitor examining the ROM. \texttt{MON E000.E00F}
runs one line without entering the prompt.}

\section{When something goes wrong}
\verb|DUMP a note| writes the whole machine state --- registers, screen, the
last 4096 instructions, your last keystrokes --- to a numbered file on the
host, for a human (or their AI) to read later. \verb|DUMP ON| does it
automatically every fifteen seconds.
